\subsection{Bayesian estimation, model selection, and prediction}
Let us consider normal linear regression models which differ in their
explanatory variables (see \citealt{Koop2003} for further details).
Suppose that we have $K$ potential explanatory variables, which means
there are $2^K$ possible models and let $M_r$ for $r=1,\dots, 2^K$
denote the $2^K$ different models under consideration.  Suppose also
that $y_i$ and $x_i$ denote the observed data on the dependent and
explanatory variables for $i=1,\ldots,N$. The observations are placed
in a $(N \times 1)$ vector $y$ and a $(N \times k_r)$ matrix $X_r$
containing the set of regressors included in model $M_r$%
\footnote{We subtract the mean from all regressors as in
	\cite{FernandezLeySteel2001a}.}.  Thus, we can write our model as
\begin{equation*}
	y=\alpha \iota_N + X_r\beta_r + \epsilon,
\end{equation*}
where $\iota_N$ is a $(N \times 1)$ vector of ones, $\beta_r$ is a
$(k_r \times 1)$ vector of unknown parameters, $\epsilon$ is a
$(N \times 1)$ vector of errors which are assumed to be normally
distributed, $\epsilon \sim N(0_N, h^{-1}I_N)$ and \textit{h} is the
error precision, which is defined as $h=\frac{1}{\sigma^2}$.
$N(\mu, \Sigma)$ denotes a normal distribution with mean $\mu$ and
variance ${\Sigma}$.  Following \cite{Koop2003}, the prior for
$\beta_r$ is normally distributed,
\begin{equation*}
	\beta_r \mid h, M_r \sim N \left(0_{k_r}, h^{-1} \left [g X_r^\top X_r  \right]^{-1}  \right),
\end{equation*}
while we use a noninformative prior for the intercept and precision,
\begin{equation*}
	p(\alpha) \propto 1, p(h)\propto \frac {1}{h}.
	\label{eq:prior}
\end{equation*}
The factor of proportionality \textit{g} is part of the so-called
Zellner's g-prior \citep{Zellner1986}. This prior is a convenient way
to specify the prior variance matrix, because it reduces the number of
prior variance parameters and considerably simplifies posterior
computations. The \pkg{gretl} package offers the five most popular
alternative Zellner's g-priors (see \citealt{FernandezLeySteel2001a}
and \citealt{Moral-Benito2010}).

\begin{itemize}\label{text:g-priors}
	\item Unit information prior (g-UIP), recommended by \cite{KassWasserman1995}:
	\begin{equation}
		g = \frac{1}{N}.
		\label{eq:g-UIP}
	\end{equation}
	\item Risk inflation criterion (g-RIC), proposed by \cite{FosterGeorge1994}:
	\begin{equation}
		g = \frac{1}{K^2}.
		\label{eq:g-RIC}
	\end{equation}
	\item Benchmark prior, recommended by \cite{FernandezLeySteel2001a}:
	\begin{equation}
		g = \left\{\begin{array}{ll}
			\frac{1}{K^2} & \mbox{for $N \le K^2$},\\
			\frac{1}{N} & \mbox{for $N>K^2$}.\end{array} \right. 
		\label{eq:Benchmark_Prior}
	\end{equation}
	\item g-HQ prior which mimics the Hannan and Quinn criterion, see \cite{FernandezLeySteel2001a}:
	\begin{equation}
		g = \frac{1}{(\text{ln}N)^3}.
		\label{eq:g-HQ}
	\end{equation}
	\item Root of g-UIP, see \cite{FernandezLeySteel2001a}:
	\begin{equation}
		g = \sqrt{\frac{1}{N}}.
		\label{eq:root-g-UIP}
	\end{equation}
\end{itemize}

By Bayes' rule, the mean of the posterior distribution of the slope
parameters $\beta_r$, conditional on model $M_r$, can be written as
\begin{equation}
	\E(\beta_r \mid y, M_r)= \left [ \left (1+g  \right ) X^\top_rX_r \right ]^{-1}X^\top_ry.
	\label{eq:posterior_mean}
\end{equation}
It is easy to see that if $g \approx 0$ the mean of the posterior
distribution (see Equation~\ref{eq:posterior_mean}) equals the OLS
estimates. The posterior variance of $\beta_r$, conditional on model
$M_r$, is given by
\begin{equation}
	\VAR \left (\beta_r \mid y, M_r \right)=\frac {Ns^2_r}{N-2}\left [ \left (1+g  \right ) X^\top_rX_r \right ]^{-1},
	\label{eq:posterior_variance2}
\end{equation}
where 
\begin{equation*}
	s^2_r=\frac{\frac{1}{1+g} y^\top P_{W_r}y+\frac{g}{1+g} \left( y-\bar y\iota_N \right)^\top \left( y-\bar y\iota_N \right)}{N}
\end{equation*}
and $P_{W_r}=I_N-W_r \left ( W^\top_rW_r \right )^{-1}W^\top_r$ for $W_r=(\iota_N,X_r)$. 

The marginal data density, conditional on model $M_r$, may be written
as
\begin{equation}
	p(y \mid M_r) \propto \left (\frac{g}{1+g}\right)^{\frac{k_r}{2}}
	\left [  \frac{1}{1+g}y^\top P_{W_r}y + \frac{g}{1+g} \left( y-\bar y\iota_N \right)^\top
	\left( y-\bar y\iota_N \right) \right]^{-\frac{N-1}{2}}.
	\label{eq:marginal_data_density}
\end{equation}

In the Bayesian approach of comparing models, it is considered useful
to employ probabilities to represent the degree of belief associated
with the alternative models. For the normal linear regression models
we can easily test two mutually exclusive (non-nested) and jointly
exhaustive models with different subset of variables.  Using Bayes'
theorem, the posterior odds ratio for a model $M_l$ against model
$M_n$ is given by
\begin{equation}
	\frac{\Prob(M_l \mid y)}{\Prob(M_n \mid y)}=\frac{\Prob(M_l)}{\Prob(M_n)}\frac{p(y \mid M_l)}{p(y \mid M_n)},
	\label{eq:posterior_odds_ratio}
\end{equation}
where $\frac{\Prob(M_l)}{\Prob(M_n)}$ is the prior odds ratio and
$\frac{p(y \mid M_l)}{p(y \mid M_n)}$ is the Bayes factor. If the
ratio $\frac{\Prob(M_l \mid y)}{\Prob(M_n \mid y)}$ is larger than 1,
we can say that the data supports model $M_l$ over model $M_n$. In our
package, we use two popular model priors.
\begin{itemize}\label{misc:priors}
	\item Binomial prior:
	\begin{equation}
		\Prob(M_r)=\theta^{k_r}(1-\theta)^{K-k_r}\; \quad\mathrm{for~} r=1, \dots, 2^K.
		\label{eq:model_prior_binomial}
	\end{equation}    
	Note that for $\theta=0.5$ we have a uniform prior on the model space, i.e., $\Prob(M_r)=2^{-K}$.
	\item Binomial-beta prior (see \citealt{GelmanCarlinSternRubin1997}):
	\begin{equation}
		\Prob(\Xi=k_r)= \frac{\Gamma (K+1)}{\Gamma(k_r+1)\Gamma(K-k_r+1)}\cdot
		\frac{\Gamma(a+k_r)\Gamma(K+b-k_r)}{\Gamma(a+b+k_r)} \cdot \frac{\Gamma(a+b)}{\Gamma(a)\Gamma(b)},
		\label{eq:model_prior_binomial-beta}
	\end{equation}
	where $\Xi$ denotes model size. 
\end{itemize}

In our package, we only need to specify the prior expected model size
$\E(\Xi) \in \left (0, K \right ]$.  Note that in case of the binomial
prior we have $\E(\Xi)=K\theta$ and the choice of $\E(\Xi)$
automatically produces a value for the prior inclusion probability
$\theta$. If we have the binomial-beta prior, the average model size
will satisfy $\E(\Xi)=\frac {a}{a+b}K$. Here, we follow
\cite{Ley_Steel_jae_2009} and fix $a=1$ hence obtaining the value of
the second hyperparameter as $b=\frac{K-\E(\Xi)}{\E(\Xi)}$.

One of the most widely accepted norms of scientific investigation is
Occam's razor.  Simply speaking, Occam's razor states that the
simplest hypothesis is also the most probable (see for example
\citealt{MadiganDavidRafteryAdrian}). In the \pkg{BMA} package we penalize
large models through the prior model probabilities
(\citealt{Osiewalski_Steel_1994}). Note that one can make $\Prob(M_r)$
a decreasing function of $k_r$ simply by defining the prior expected
model size $\E(\Xi)$. The smaller the prior expected model size
$\E(\Xi)$, the less probable are large models.

It is easy to show that the posterior probability of model $M_l$ is given by
\begin{equation*}
	\Prob(M_l \mid y)= \frac{\Prob(M_l)p(y \mid M_l)}{\sum_{r=1}^{2^K} \Prob(M_r)p(y \mid M_r)}.
\end{equation*}

The posterior density of vector $\beta$ is the average
of the posterior densities $p(\beta_r \mid y, M_r)$ conditional on the models
\begin{equation*}
	p(\beta \mid y)=\sum_{r=1}^{2^K} \Prob(M_r \mid y)p(\beta_r \mid y, M_r).
\end{equation*}

Once the model posterior probabilities have been calculated, we can also easily
evaluate the mean and variance of the posterior distribution of slope parameters
(see \citealt{Leamer1978}).
\begin{equation}
	\E(\beta \mid y)=\sum_{r=1}^{2^K} \Prob(M_r \mid y)\E(\beta_r \mid y, M_r)
	\label{eq:mean_averaging}
\end{equation} 
and
\begin{multline}
	\VAR(\beta \mid y) = \\\sum_{r=1}^{2^K}\Prob(M_r \mid y)\VAR(\beta_r \mid y, M_r)+
	\sum_{r=1}^{2^K} \Prob(M_r \mid y) \left (\E(\beta_r \mid y, M_r)-\E(\beta \mid y) \right)^2
	\label{eq:variance_averaging}
\end{multline}

In a similar manner, we can find other characteristics of the
posterior distribution (see for example \citealt[p.~266]{Koop2003}).
We might be also interested in the estimates of the posterior
inclusion probability $\Prob(i \mid y)$ (PIP), i.e., the probability
that, conditional on the data, but unconditional with respect to a
specific model, the variable $x_i$ is relevant in explaining the
dependent variable \textit{y} (see \citealt{Leamer1978,
	Mitchell_Beauchamp_1988,DoppelhoferWeeks2009}).  The PIP is
calculated as the sum of the posterior model probabilities for all of
the models including variable $x_i$.

In Bayesian model averaging approach, predictions of the dependent
variable are made by taking a weighted average of the individual model
forecasts. As described in \cite{FernandezLeySteel2001b}, prediction
of an unobserved data point $y_f$ is based on calculating
\begin{multline}
	p(y_f\mid y) = \sum_{r=1}^{2^K} f_S ( y_f \mid N-1,\bar y + \frac{1}{g+1}x^\top_{f,r}\beta^*_r,\\
	\frac{N-1}{Ns^2_r} \{1+\frac{1}{N}+\frac{1}{g+1}x^\top_{f,r}(X^\top_jX_j)^{-1}x_{f,r}\}^{-1})\Prob(M_r\mid y),    
	\label{eq:predictive_pdf}
\end{multline}
where $f_S ( y_f \mid v, b, a)$ denotes the univariate Student-$t$
density with $v$ degrees of freedom, mean $b$ and precision $a$ (with
variance $ v/ \{a(v-2)\}$).  In addition,
$\beta^*_r=\left [X^\top_rX_r \right ]^{-1}X^\top_ry$ is the ordinary
least squares estimator for $\beta_r$ and $x_{f,r}$ groups the $r$
elements of $x_{f}$ corresponding to the regressors in model $M_r$. 
